
@@M0@@
donne par adjonction un homomorphisme~:
@@M1@@
et on vérifie immédiatement sur chaque fibre que c'est un isomorphisme.

2.3.2. Soient maintenant @@M2@@ une @@M3@@-Algèbre graduée quasi-cohérente, @@M4@@ un @@M5@@-Module gradué, et @@M6@@. Pour tout ouvert affine @@M7@@ de @@M8@@, et toute section @@M9@@ de @@M10@@ sur @@M11@@, homogène de degré @@M12@@, on note @@M13@@ la composante homogène de degré @@M14@@ du Module de fraction @@M15@@; on a donc :
@@M16@@
Si @@M17@@ est un ouvert affine de @@M18@@, alors @@M19@@ est la restriction à @@M20@@ de @@M21@@.

On considère la base d'ouverts de @@M22@@ formée des ouverts affines :
@@M23@@
On note @@M24@@ le morphisme structural @@M25@@, et @@M26@@ le morphisme affine :
@@M27@@
On pose alors :
@@M28@@
avec les notations de 2.3.1. Nous allons montrer que les @@M29@@ se recollent quand @@M30@@ et @@M31@@ varient, définissant ainsi un @@M32@@-Module noté @@M33@@. Le recollement pour @@M34@@ variable et @@M35@@ fixe étant évident, on est ramené (cf. EGA II 2.5.2) à montrer que si @@M36@@ est une section de @@M37@@ sur @@M38@@, homogène de degré @@M39@@, alors la restriction de @@M40@@ à @@M41@@ est canoniquement isomorphe à @@M42@@. On considère :
